\subsection{Primary Decomposition}

\begin{theorem} \label{thm:complex-primary-decomposition-generalized-eigenspaces} \hlt{Primary Decomposition over $\mathbb{C}$}\\
Let $\mathcal{V}\neq\{\mathbf{0}\}$ be finite-dim over $\mathbb{C}$, let $T\in\LL(\mathcal{V})$, and let $\lambda_1,\ldots,\lambda_m$ be the distinct eigenvalues of $T$. Then
\begin{enumerate}[label=(\roman*)]
\setlength{\itemsep}{0pt}
\item $\mathcal{V}=\mathcal{G}_{\lambda_1}(T)\oplus\cdots\oplus \mathcal{G}_{\lambda_m}(T)$;
\item each $(T-\lambda_jI_{\mathcal{V}})|_{\mathcal{G}_{\lambda_j}(T)}$ is nilpotent.
\end{enumerate}
\begin{proof}
Let $n=\dim(\mathcal{V})$.
\begin{enumerate}[label=(\roman*)]
\setlength{\itemsep}{0pt}
\item
Proceed by induction on $n$. If $n=1$, then $T$ has single eigenvalue $\lambda_1$ and $\mathcal{V}=\mathcal{G}_{\lambda_1}(T)$.\\
Assume $n>1$. Let $S=T-\lambda_1I_{\mathcal{V}}$, $\mathcal{U}\defeq\operatorname{range}(S^n)$.\\
By Theorems \ref{thm:kernel-range-stable-direct-sum} and \ref{thm:generalized-eigenspace-stabilized-kernel}, $\mathcal{V} =\ker(S^n)\oplus\operatorname{range}(S^n) =\mathcal{G}_{\lambda_1}(T)\oplus\mathcal{U}$.
The subspace $\mathcal{U}$ is $T$-invariant by Proposition \ref{prop:kernel-range-polynomial-in-operator-invariant}, and $\dim(\mathcal{U})<n$ because $\mathcal{G}_{\lambda_1}(T)\neq\{\mathbf{0}\}$. If $\mathcal{U}=\{\mathbf{0}\}$, then (i) follows.\\
Otherwise distinct eigenvalues of $T|_{\mathcal{U}}$ are $\lambda_2,\ldots,\lambda_m$: every eigenvalue of $T|_{\mathcal{U}}$ is an eigenvalue of $T$ but not $\lambda_1$, while an eigenvector $\mathbf{v}$ of $T$ for $\lambda_j$, $j\geq2$, satisfies $S^n\mathbf{v}=(\lambda_j-\lambda_1)^n\mathbf{v}$ and hence lies in $\mathcal{U}$.\\
Induction gives $\mathcal{U} =\mathcal{G}_{\lambda_2}(T|_{\mathcal{U}})\oplus\cdots\oplus \mathcal{G}_{\lambda_m}(T|_{\mathcal{U}})$. For $j\geq2$, $\mathcal{G}_{\lambda_j}(T|_{\mathcal{U}})=\mathcal{G}_{\lambda_j}(T)$.\\
The inclusion $\subseteq$ follows from the definition. Conversely, if $\mathbf{v}\in\mathcal{G}_{\lambda_j}(T)$, write $\mathbf{v}=\mathbf{g}+\mathbf{u}$ with $\mathbf{g}\in\mathcal{G}_{\lambda_1}(T)$ and $\mathbf{u}\in\mathcal{U}$. For some $r$, $(T-\lambda_jI_{\mathcal{V}})^r\mathbf{v}=\mathbf{0}$. Since both summands are $T$-invariant, the two terms $(T-\lambda_jI_{\mathcal{V}})^r\mathbf{g}$ and $(T-\lambda_jI_{\mathcal{V}})^r\mathbf{u}$ remain in the two direct summands and have sum zero, so both are zero.\\
Thus $\mathbf{g}\in\mathcal{G}_{\lambda_1}(T)\cap\mathcal{G}_{\lambda_j}(T)=\{\mathbf{0}\}$ by Theorem \ref{thm:distinct-generalized-eigenvectors-independent}, and $\mathbf{v}=\mathbf{u}\in\mathcal{G}_{\lambda_j}(T|_{\mathcal{U}})$. Hence (i) follows.
\item
By Proposition \ref{prop:generalized-eigenspaces-invariant}, the restriction is an operator on $\mathcal{G}_{\lambda_j}(T)$. By Theorem \ref{thm:generalized-eigenspace-stabilized-kernel}, $(T-\lambda_jI_{\mathcal{V}})^n\mathbf{v}=\mathbf{0} \ \forall \mathbf{v}\in\mathcal{G}_{\lambda_j}(T)$. Hence $n$th power of $(T-\lambda_jI_{\mathcal{V}})|_{\mathcal{G}_{\lambda_j}(T)}$ is zero.
\end{enumerate}
\hlt{Structural Commentary}: Peel off one generalized eigenspace, then induct on the invariant range.\\
\hlt{Key Insight}: Stable kernel-range splitting isolates one eigenvalue at a time.
\end{proof}
\end{theorem}

\begin{figure}[H]
  \centering
  \begin{tikzpicture}[
  >=Stealth,
  every node/.style={font=\small},
  block/.style={draw=blue!55!black, fill=blue!6, rounded corners=3pt, minimum width=2.1cm, minimum height=1.05cm, align=center},
  arrow/.style={->, line width=0.85pt, blue!70!black},
  nil/.style={->, line width=0.75pt, red!65!black}
]
  \node[block] (G1) at (0,0) {$G(\lambda_1,T)$\\$T=\lambda_1 I+N_1$};
  \node[block] (G2) at (2.95,0) {$G(\lambda_2,T)$\\$T=\lambda_2 I+N_2$};
  \node[block] (Gm) at (6.16,0) {$G(\lambda_m,T)$\\$T=\lambda_m I+N_m$};

  \node at (4.58,0) {$\cdots$};
  \node[above=0.62cm of G2, font=\small] (V) {$\mathcal{V}=G(\lambda_1,T)\oplus G(\lambda_2,T)\oplus\cdots\oplus G(\lambda_m,T)$};

  \draw[arrow] (V.south west) to[out=-110,in=90] (G1.north);
  \draw[arrow] (V.south) -- (G2.north);
  \draw[arrow] (V.south east) to[out=-70,in=90] (Gm.north);

  \draw[nil] ($(G1.south west)+(0.25,0.18)$) .. controls +(0,-0.55) and +(0,-0.55) .. ($(G1.south east)+(-0.25,0.18)$)
    node[midway, below=5pt] {$N_1^{r_1}=0$};
  \draw[nil] ($(G2.south west)+(0.25,0.18)$) .. controls +(0,-0.55) and +(0,-0.55) .. ($(G2.south east)+(-0.25,0.18)$)
    node[midway, below=5pt] {$N_2^{r_2}=0$};
  \draw[nil] ($(Gm.south west)+(0.25,0.18)$) .. controls +(0,-0.55) and +(0,-0.55) .. ($(Gm.south east)+(-0.25,0.18)$)
    node[midway, below=5pt] {$N_m^{r_m}=0$};
\end{tikzpicture}

  \caption{Primary decomposition: on $\mathcal{G}_{\lambda_j}(T)$, $T=\lambda_j I$ plus nilpotent.}
  \label{fig:primary-decomposition-generalized-eigenspaces}
\end{figure}

\begin{corollary} \label{cor:complex-basis-of-generalized-eigenvectors} \hlt{Basis of Generalized Eigenvectors}\\
Let $\mathcal{V}\neq\{\mathbf{0}\}$ be finite-dimensional over $\mathbb{C}$ and $T\in\LL(\mathcal{V})$. Then $\mathcal{V}$ has a generalized-eigenvector basis for $T$.
\begin{proof}
Let $\lambda_1,\ldots,\lambda_m$ be the distinct eigenvalues of $T$. By Theorem \ref{thm:complex-primary-decomposition-generalized-eigenspaces}, $\mathcal{V}=\mathcal{G}_{\lambda_1}(T)\oplus\cdots\oplus \mathcal{G}_{\lambda_m}(T)$. Choose a basis of each generalized eigenspace and concatenate the bases. Directness makes the concatenated list a basis of $\mathcal{V}$, and every nonzero vector in each list is a generalized eigenvector.\\
\hlt{Structural Commentary}: A direct sum basis is obtained by choosing bases in each summand.\\
\hlt{Key Insight}: Bases of the direct-sum pieces concatenate.
\end{proof}
\end{corollary}

\begin{definition} \label{def:eigenvalue-multiplicity-generalized-eigenspace} \hlt{Multiplicity of an Eigenvalue}\\
Let $\mathcal{V}$ be finite-dimensional and let $T\in\LL(\mathcal{V})$.\\
The \hlt{multiplicity} of an eigenvalue $\lambda$ of $T$ is $\dim \mathcal{G}_{\lambda}(T)=\dim\ker\bigl((T-\lambda I_{\mathcal{V}})^{\dim(\mathcal{V})}\bigr)$.\\
This is the algebraic multiplicity in generalized-eigenspace language.
\end{definition}

\begin{proposition} \label{prop:sum-of-complex-eigenvalue-multiplicities} \hlt{Sum of Multiplicities}\\
Let $\mathcal{V}\neq\{\mathbf{0}\}$ be a finite-dimensional complex vector space and let $T\in\LL(\mathcal{V})$. Then the sum of the multiplicities of all eigenvalues of $T$ equals $\dim(\mathcal{V})$.
\begin{proof}
Let $\lambda_1,\ldots,\lambda_m$ be the distinct eigenvalues of $T$. By Theorem \ref{thm:complex-primary-decomposition-generalized-eigenspaces}, $\mathcal{V}=\mathcal{G}_{\lambda_1}(T)\oplus\cdots\oplus \mathcal{G}_{\lambda_m}(T)$.\\
Taking dimensions, by Corollary \ref{cor:dimension-direct-sum}, $\dim(\mathcal{V})=\sum\limits_{j=1}^{m}\dim \mathcal{G}_{\lambda_j}(T)$, exactly the sum of the multiplicities.\\
\hlt{Structural Commentary}: Multiplicity is dimension of the generalized-eigenspace summand.\\
\hlt{Key Insight}: Multiplicities are dimensions of the direct-sum pieces.
\end{proof}
\end{proposition}

\begin{theorem} \label{thm:complex-block-diagonal-generalized-eigenbasis} \hlt{Block Diagonal Form with Upper-Triangular Primary Blocks}\\
Let $\mathcal{V}\neq\{\mathbf{0}\}$ be finite-dimensional over $\mathbb{C}$, and let $\lambda_1,\ldots,\lambda_m$ be the distinct eigenvalues of $T\in\LL(\mathcal{V})$, with multiplicities $d_1,\ldots,d_m$. Then $T$ has a matrix of the form
\begin{equation*}
\begin{bmatrix}
A_1 & & 0\\
& \ddots & \\
0 & & A_m
\end{bmatrix},
\end{equation*}
where each $A_j\in\operatorname{Mat}(d_j,d_j,\mathbb{C})$ is upper triangular with diagonal $\lambda_j,\ldots,\lambda_j$.
\begin{proof}
Set $\mathcal{U}_j=\mathcal{G}_{\lambda_j}(T)$ and $N_j=(T-\lambda_jI_{\mathcal{V}})|_{\mathcal{U}_j}$.\\
By Theorem \ref{thm:complex-primary-decomposition-generalized-eigenspaces}, the $\mathcal{U}_j$ form a $T$-invariant direct-sum decomposition and each $N_j$ is nilpotent.\\
Choose a basis of each $\mathcal{U}_j$ in which $N_j$ is strictly upper triangular. Then $T|_{\mathcal{U}_j}=\lambda_jI_{\mathcal{U}_j}+N_j$ has an upper triangular matrix with diagonal $\lambda_j,\ldots,\lambda_j$. Concatenating these bases gives the desired block diagonal matrix.\\
\hlt{Structural Commentary}: Primary decomposition gives the blocks; nilpotent triangularisation shapes them. \\
\hlt{Key Insight}: On each generalized eigenspace, $T$ is eigenvalue plus nilpotent.
\end{proof}
\end{theorem}

\begin{figure}[H]
  \centering
  \begin{tikzpicture}[
  every node/.style={font=\small},
  mat/.style={draw=gray!60, fill=gray!5, rounded corners=2pt, inner xsep=5pt, inner ysep=4.5pt, align=center},
  block/.style={draw=blue!60!black, fill=blue!7, minimum width=1.48cm, minimum height=0.53cm, align=center},
  arrow/.style={->, line width=0.85pt, blue!70!black}
]
  \node[mat] (M) at (0,0) {
    $\begin{bmatrix}
      A_1 & & 0\\
      & \ddots & \\
      0 & & A_m
    \end{bmatrix}$};

  \node[block] (A1) at (3.26,0.72) {$A_1$};
  \node[block] (Aj) at (3.26,0) {$A_j$};
  \node[block] (Am) at (3.26,-0.72) {$A_m$};

  \draw[arrow] (M.east) -- (A1.west);
  \draw[arrow] (M.east) -- (Aj.west);
  \draw[arrow] (M.east) -- (Am.west);
\end{tikzpicture}

  \caption{Primary blocks: $A_j=\lambda_jI+N_j$, with $N_j$ strictly upper triangular.}
  \label{fig:primary-block-diagonal-form}
\end{figure}
